\chapter{Objetivos do Curso}

O \ppccurso destina-se a preparar tecnólogos(as) para atuar de forma integrada nas áreas de instrumentação, controle, automação industrial, redes industriais e, como diferencial deste curso, robótica e inteligência artificial aplicada à automação. O profissional formado pela FEELT/UFU será capaz de agir com base em conhecimentos técnico-aplicados sólidos e manter postura crítica, ética e alinhada às demandas do setor produtivo. Espera-se que demonstre iniciativa, competência técnica, responsabilidade profissional e compromisso com a segurança e a sustentabilidade dos processos industriais.

\section{Objetivo Geral}

Formar tecnólogos(as) em Automação Industrial com domínio técnico-aplicado das tecnologias de instrumentação, controle, automação, redes industriais, acionamentos e, na ênfase deste curso, robótica e sistemas inteligentes, capacitados(as) a projetar, implementar, operar e manter sistemas automatizados industriais e a atuar de forma crítica e aplicada na identificação e resolução de problemas do setor produtivo. O(a) egresso(a) deverá considerar os aspectos técnicos, econômicos, sociais, ambientais e de segurança do trabalho inerentes à prática profissional, mantendo uma postura ética em consonância com as transformações da Indústria 4.0 e as demandas do mercado de trabalho brasileiro.

\section{Objetivos Específicos}

Considerando as Diretrizes Curriculares Nacionais Gerais para a Educação Profissional e Tecnológica~\cite{cne_cp_1_2021} e o Catálogo Nacional de Cursos Superiores de Tecnologia~\cite{mec_cncst_2024}, o \ppccurso busca:

\begin{itemize}
    \item Assegurar formação técnico-aplicada sólida nas tecnologias centrais da automação industrial — eletricidade, eletrônica, instrumentação, controle, CLP/SCADA e redes industriais — diretamente orientada ao exercício profissional;
    \item Promover a integração curricular entre os núcleos de formação básica, tecnológica/profissional, humanística e de integração/extensão, assegurando progressão adequada e interdisciplinaridade;
    \item Capacitar os(as) futuros(as) tecnólogos(as) em Automação Industrial a incorporar princípios de sustentabilidade, eficiência energética e segurança do trabalho em sua prática profissional;
    \item Garantir atualização permanente do currículo, flexibilizando os conteúdos mais sujeitos à obsolescência tecnológica por meio de disciplinas optativas;
    \item Favorecer a realização de projetos integradores e atividades curriculares de extensão, visando o diálogo entre diferentes áreas de conhecimento e a comunidade;
    \item Integrar teoria e prática, conectando os conteúdos curriculares às realidades da atuação profissional, por meio de aulas práticas, projetos, estudos de caso e experiências aplicadas em laboratório;
    \item Reconhecer e incorporar ao percurso formativo as competências desenvolvidas em contextos não formais de aprendizagem (estágio, monitorias, extensão, vivência profissional);
    \item Acompanhar as transformações do mercado de trabalho e da indústria, fomentando respostas éticas, criativas e socialmente responsáveis às demandas consolidadas e emergentes da automação industrial;
    \item Assegurar a formação de profissionais capazes de:
    \begin{itemize}
        \item Compreender e atuar sobre questões técnicas, éticas, legais e de segurança do trabalho que permeiam a prática profissional em ambientes industriais automatizados;
        \item Atender de forma direta e aplicada às necessidades do setor produtivo por meio da automação de processos e sistemas;
        \item Contribuir com visão técnica e aplicada para a modernização de plantas e processos industriais;
        \item Cooperar em equipes multidisciplinares de engenharia, operação e manutenção, atuando com senso de responsabilidade;
        \item Utilizar de forma racional os recursos disponíveis, com foco na sustentabilidade, na segurança e na eficiência dos processos automatizados.
    \end{itemize}
\end{itemize}
